%%%%%%%%%%%%%%%%%%%%%%%%%%%%%%%%%%%%%%%%%%%%%%%%%%%%%%%%%%%%%%%%%%%%%%%
%% Optionen zum Layout des Artikels                                  %%
%%%%%%%%%%%%%%%%%%%%%%%%%%%%%%%%%%%%%%%%%%%%%%%%%%%%%%%%%%%%%%%%%%%%%%%
\documentclass[%
a4paper,							% alle weiteren Papierformat einstellbar
oneside,							% einseitiger Druck
%twoside,							% Doppelseiten
%twocolumn,						% zweispaltiger Satz
%landscape,						% Querformat
12pt,								% Schriftgröße (12pt, 11pt (Standard))
%BCOR1cm,							% Bindekorrektur, bspw. 1 cm
%DIVcalc,							% führt die Satzspiegelberechnung neu aus
%halfparskip*,				% Absatzformatierung s. scrguide 3.1
headsepline,					% Trennline zum Seitenkopf	
footsepline,					% Trennline zum Seitenfuß
titlepage,						% Titelei auf eigener Seite
%normalheadings,			% Überschriften etwas kleiner (smallheadings)
%leqno,   						% Nummerierung von Gleichungen links
%fleqn,								% Ausgabe von Gleichungen linksbündig
%draft								% überlangen Zeilen in Ausgabe gekennzeichnet
toc=listof,						% Abb.-Verzeichnis im Inhalt
toc=bibliography			% Tab.-Verzeichnis im Innhalt
]
{scrbook}


%% Deutsche Anpassungen %%%%%%%%%%%%%%%%%%%%%%%%%%%%%%%%%%%%%
\usepackage[ngerman]{babel}
\usepackage[ngerman]{translator}
\usepackage[T1]{fontenc}
\usepackage[utf8]{inputenc} %UTF-8 als Charset
\usepackage{lmodern} % Type1-Schriftart für nicht-englische Texte
\usepackage[autostyle=true,german=quotes]{csquotes}





%% Globale Variablen und Funktionen %%%%%%%%%%%%%%%%%%%%%%%%%%%%%%%%%%%%%
\renewcommand{\author}{Autor}
\renewcommand{\title}{Ein geeigneter Titel}
\newcommand{\uni}{Duale Hochschule Baden-Württemberg Mannheim}
\newcommand{\art}{Bachelorarbeit}
\newcommand{\studiengang}{International Management for Business and Information Technology}
\newcommand{\ort}{Mannheim}
\newcommand{\abgabedatum}{xx.xx.xx}
\newcommand{\kurs}{Kurs}
\newcommand{\matrikelnummer}{M-Nr.}
\newcommand{\firma}{Unternehmen}
\newcommand{\firmenort}{Firmenort}
\newcommand{\firmenbetreuer}{Betreuer}
\newcommand{\hochschulbetreuer}{Betreuer}


% freie Seiten ohne Seitencounter zu erhöhen
\newcommand\blankpage{%
	\newpage
	\thispagestyle{empty}
	\mbox{}
	\addtocounter{page}{-1}
	\newpage
}

% Abstractumgebung definieren
\newenvironment{abstract}%
    {\thispagestyle{empty}\null\vfill\begin{center}%
    \bfseries\abstractname\end{center}}%
    {\vfill\null}

%% allgemeine Paketimports und Einstellungen %%%%%%%%%%%%%%%%%%%%%%%%%%%%%%%%%%%%%


\usepackage{color}
\definecolor{LinkColor}{rgb}{0,0,0}
\definecolor{ListingBackground}{rgb}{0.92,0.92,0.92}

% Verlinkunggen innheralb der PDF, PDF Attribute, etc. 
\usepackage[%
	pdftitle={\title},
	pdfauthor={\author},
	pdfsubject={\art},
	pdfcreator={pdflatex, LaTeX with KOMA-Script},
	pdfpagemode=UseOutlines, % Beim Oeffnen Inhaltsverzeichnis anzeigen
	pdfdisplaydoctitle=true, % Dokumenttitel statt Dateiname anzeigen.
	pdflang=de % Sprache des Dokuments.
]{hyperref} 

% (Farb-)einstellungen für die Links im PDF
\hypersetup{%
	colorlinks=true, % Aktivieren von farbigen Links im Dokument
	linkcolor=LinkColor, % Farbe festlegen
	citecolor=LinkColor,
	filecolor=LinkColor,
	menucolor=LinkColor,
	urlcolor=LinkColor,
	bookmarksnumbered=false % Überschriftsnummerierung im PDF Inhalt anzeigen.
}

\usepackage{bookmark}
\bookmarksetup{
  numbered, 
  open,
}

% Einstellungen für das Inhaltsverzeichnis
\usepackage{tocstyle}
\newtocstyle[KOMAlike][leaders]{alldotted}{}
\usetocstyle{alldotted}

% Glossar und Abkürzungsverzeichnis, Eintrag in TOC, keine Seitenzahlen hinter Eintrag
\usepackage[toc, acronym, nonumberlist]{glossaries}
\makeindex
\makeglossaries
\loadglsentries {content/glossar}
\loadglsentries {content/acronyme}



\usepackage[onehalfspacing]{setspace} %Zeilenabstand setzen

\usepackage{geometry} %Seitenränder
\geometry{a4paper, left=30mm, right=25mm,}

\usepackage{graphicx} % Zum Laden von Grafiken

\usepackage[%
	style=authoryear,	%Harvard Style
	backend=biber,			%Biber backend für UTF8 Support
	isbn=false,
	doi=false
]{biblatex} % Bibliographiestil und Package
\bibliography{content/bibliographie}	% Literaturverzeichnis


\clubpenalty = 10000 
\widowpenalty = 10000 
\displaywidowpenalty = 10000

\begin{document}

\pagestyle{empty} %%Keine Kopf-/Fusszeilen auf den ersten Seiten.


%%%%%%%%%%%%%%%%%%%%%%%%%%%%%%%%%%%%%%%%%%%%%%%%%%%%%%%%%%%%%%%%%%%%%%%
%% Ihr Artikel                                                       %%
%%%%%%%%%%%%%%%%%%%%%%%%%%%%%%%%%%%%%%%%%%%%%%%%%%%%%%%%%%%%%%%%%%%%%%%

\pdfbookmark[section]{Titelseite}{title}
\begin{titlepage}

\newgeometry{
	top=3cm,
	bottom=2cm
}

\newcommand{\HRule}{\rule{\linewidth}{0.5mm}} % Defines a new command for the horizontal lines, change thickness here

 
%----------------------------------------------------------------------------------------
%	HEADING SECTIONS
%----------------------------------------------------------------------------------------

\begin{minipage}[c]{0.5\textwidth}
\includegraphics[height=3.2cm]{images/dhbw.png}
\end{minipage}
\hspace{.4cm}
\begin{minipage}[c]{0.5\textwidth}
\includegraphics[height=3.2cm]{images/firma.png}
\end{minipage}

\center % Center everything on the page

\vspace{1cm}
\textsc{\LARGE \uni}\\[.5cm] % Name of your university/college
%----------------------------------------------------------------------------------------
%	TITLE SECTION
%----------------------------------------------------------------------------------------

\HRule \\[.4cm]
{ \huge \bfseries \title}\\ % Title of your document
\HRule \\[.4cm]
\vspace{.5cm}
\textsc{\Large \art}\\[.5cm]
\textit{im Studiengang}\\
\textsc{{\Large \studiengang}}\\[0.5cm] % Major heading such as course name
\textit{von}\\
\textsc{\Large \author}\\[2cm]
\normalsize
%----------------------------------------------------------------------------------------
%	AUTHOR SECTION
%----------------------------------------------------------------------------------------
\begin{tabbing}
mmmmmmmmmmmmmmmmmmmmmmmmmm     \= \kill
		\textbf{Matrikelnummer, Kurs}			 \> \matrikelnummer, \kurs\\
		\textbf{Firma}      \> \firma, \firmenort\\
		\textbf{Firmenbetreuer}              \> \firmenbetreuer\\
		\textbf{Hochschulbetreuer}             \>  \hochschulbetreuer
\end{tabbing}
%----------------------------------------------------------------------------------------
%	DATE SECTION
%----------------------------------------------------------------------------------------
\vspace{1cm}
{\large \abgabedatum}\\ % Date, change the \today to a set date if you want to be precise


\vfill % Fill the rest of the page with whitespace

\end{titlepage}
\restoregeometry

% Pagenumbering auf römisch für die Preambel
\renewcommand{\thepage}{\Roman{page}}
\setcounter{page}{1}

%% Ehrenwörtliche Erklärung %%%%%%%%%%%%%%%%%%%%%%%%%%%%%%%%%%%%%%%
\phantomsection
\addcontentsline{toc}{chapter}{Ehrenwörtliche Erklärung}
\input{content/ehrenwörtliche}
\newpage


%% Abstract %%%%%%%%%%%%%%%%%%%%%%%%%%%%%%%%%%%%%%%
\phantomsection
\renewcommand\abstractname{Abstract} % Überschrift des Abstracts
\begin{abstract}
	\input{content/abstract}
\end{abstract}
\addcontentsline{toc}{chapter}{Abstract}


%% Inhaltsverzeichnis %%%%%%%%%%%%%%%%%%%%%%%%%%%%%%%%%%%%%%%
\phantomsection
\tableofcontents			% Inhaltsverzeichnis
\addcontentsline{toc}{chapter}{Inhaltsverzeichnis}


% Pagenumbering auf arabisch für den Hauptteil
\newpage
\renewcommand{\thepage}{\arabic{page}}
\setcounter{page}{1}

%\pagestyle{empty}		% keine Kopf und Fußzeile (k. Seitenzahl)
\pagestyle{headings}	% lebender Kolumnentitel  


%% eigentlicher Inhalt %%%%%%%%%%%%%%%%%%%%%%%%%%%%%%%%%%%%%%%
\input{content/einführung}



%% Appendix %%%%%%%%%%%%%%%%%%%%%%%%%%%%%%%%%%%%%%%

% Pagenumbering wieder auf römisch für Appendix
\newpage
\renewcommand{\thepage}{\Roman{page}}
\setcounter{page}{4}

%% Verzeichnisse %%%%%%%%%%%%%%%%%%%%%%%%%%
\listoftables				% Tabellenverzeichnis
\listoffigures				% Abbildungsverzeichnis

% getrennte Quellenverzeichnisse
% für ein einzelnes Quellenverzeichnis nur \printbibliography ohne Optionen verwenden
% Filter für die "`Hauptbibliografie"' bauen
% !Die in diesem Filter genannten types unten als 'nottype' aufführen
\defbibfilter{main}{
	type=article
	or type=book
}

\printbibliography[filter=main] % Hauptbibliografie drucken
\printbibliography[title={Sonstige Quellen}, nottype=book, nottype=article] % Sonstige Quellen drucken


\printglossary[type=\acronymtype,title=Abkürzungsverzeichnis]
\printglossary[type=main,title=Glossar]

\end{document}